%% ma: L10-B14-0033
%% lop: 10
%% chuong: 5
%% bai: 14
%% dang: 10.14.5
%% loai: DS
%% muc: [NB, TH, NB, VD]
%% dap_an: "ĐSSĐ"
%% nguon: "chuyen-de-cac-so-dac-trung-cua-mau-so-lieu-khong-ghep-nhom-toan-10-chuong-trinh-moi.pdf (GV. Phan Nhật Linh) (P04), Bài 03 Dạng 1 Phần II Câu 10"
%% kien_thuc: "Lớp 10 Bài 13 (số trung bình, tứ phân vị), Bài 14 (khoảng biến thiên, khoảng tứ phân vị, giá trị bất thường)"
%% trang_thai: chinh-thuc
%% ghi_chu: "Cùng số liệu Bài tập tự luận 6 nhưng khác loại câu và câu hỏi. Tự duyệt theo yêu cầu giáo viên 22/10/2026."
\begin{ex}
Điểm khảo sát chất lượng học sinh giỏi của $15$ học sinh khối 10 được thống kê như sau:
\begin{center}
\begin{tabular}{|c|c|c|c|c|c|c|c|c|c|c|c|c|c|c|}\hline
$5$ & $9$ & $6$ & $8$ & $6$ & $6$ & $7$ & $4$ & $6$ & $7$ & $7$ & $9$ & $8$ & $5$ & $8$\\\hline
\end{tabular}
\end{center}
Xét tính đúng sai của các khẳng định sau:
\choiceTF{\True Khoảng biến thiên của mẫu số liệu là $5$}{Khoảng tứ phân vị của mẫu số liệu là $1$}{Trung vị của mẫu số liệu là $4$}{\True Mẫu số liệu không có giá trị bất thường}
\loigiai{Sắp xếp: $4$; $5$; $5$; $6$; $6$; $6$; $6$; $7$; $7$; $7$; $8$; $8$; $8$; $9$; $9$.

a) Đúng: $R=9-4=5$.

b) Sai: $Q_1=6$, $Q_3=8$, $\Delta_Q=2$.

c) Sai: $M_e=x_8=7$.

d) Đúng: $Q_1-1{,}5\Delta_Q=3$; $Q_3+1{,}5\Delta_Q=11$; mọi giá trị đều thuộc $[3;11]$.}
\end{ex}
